\documentclass[compress]{beamer} %en animé
%\documentclass[trans]{beamer} %en superposé

%%%%%%%%%%%%%% LES PACKAGES
%-------------------
\usepackage[utf8]{inputenc}
\usepackage[T1]{fontenc}
\usepackage[frenchb]{babel}
%\usepackage[english]{babel}

%-------------------
\usepackage{minted} %pour inclure du code
\setminted[ocaml]{	                               
	%bgcolor = black!3!white,                           %--fond 
	frame = leftline,	framesep = 6pt, rulecolor= expli, %-- cadre
	linenos=true, numbersep=4pt, xleftmargin=20pt,      %--numéro de ligne
	breaklines=true,                                    %--découpage des longues lignes
	tabsize=2,     %--tabulations    
	fontsize=\scriptsize,
}

\setminted[c]{	                               
	%bgcolor = black!3!white,                           %--fond 
	frame = leftline,	framesep = 6pt, rulecolor= expli, %-- cadre
	linenos=true, numbersep=4pt, xleftmargin=20pt,      %--numéro de ligne
	breaklines=true,                                    %--découpage des longues lignes
	tabsize=2,                                          %--tabulations
}

%-------------------
\usepackage[ruled,vlined]{algorithm2e} %pour le pseudo-code
%ajouter l'option linesnumbered si on veut les numéros de lignes
\SetAlgorithmName{Algorithme}{Algo}{Liste des algorithmes}
\SetFuncArgSty{textup}%style des arguments des fonctions pas en gras
\SetArgSty{textup}%style des condidtions des if, for, while pas en gras


\SetKwFor{ForEach}{pour tout}{faire}{}
\SetKwFor{For}{pour}{faire}{finpour}
\SetKwIF{Si}{SinonSi}{Sinon}{si}{alors}{sinon si}{sinon}{}
\SetKwInput{Input}{Entrée}
\SetKwInput{Output}{Sortie}
\SetKwProg{myproc}{Procedure}{:}{}
\SetKw{Return}{retourner}
\SetKwComment{tcc}{(*}{*)}
\SetKwFor{Tq}{tant que}{faire}{}
\SetKwRepeat{Repeter}{répéter}{jusqu’à}

%-------------------
\usepackage{amsmath,amsfonts,amssymb}
\usepackage{textcomp,lmodern}%pour l'euro
\usepackage{cancel}%pour barrer

%-------------------
\usepackage{array,multirow}
\usepackage{tabularx}%tableaux élastiques


%-------------------
\usepackage{tikz} %graphiques et dessins
\usetikzlibrary{shapes}%pr écrire \node[ellipse] par ex
\usetikzlibrary{positioning} %pr écrire \node[below rigth = 3pt and 5pt]
\usetikzlibrary{decorations.pathreplacing}%pr les accolades
\usetikzlibrary{patterns}%pour les hachures

%-------------------
\usepackage{caption}%les personaliser (les centrer par ex)
\usepackage{changepage}%pour élargir la page avec adjustwidth
\usepackage{calc}%pour pouvoir écrire \textwidth-1cm


%-------------------
\usepackage{xspace}%espaces après les commandes texte


%------------------

%-------------------



\usepackage[clock]{ifsym}

%%%%%%%%%%%%%% THEME BEAMER
%votre nom court qui apparaîtra dans l'en tête
\newcommand{\initiales}{Yann Miquel--Erdmann }
\usetheme{Montpellier}
\usecolortheme{seahorse}

%--------marges
\setbeamersize{text margin left= 0.7cm}
\setbeamersize{text margin right= 0.7cm}

%--------tête et pieds
\setbeamertemplate{navigation symbols}{}
\setbeamertemplate{footline}[frame number]
\setbeamertemplate{headline}{
  %la premiere ligne
  	\begin{beamercolorbox}[ht=0.42cm, vmode]{section in head/foot}
	\hspace{0.4cm} \insertshorttitle 
	\hspace*{0.1cm}- \initiales - {\insertshortdate}
	\vspace*{0.08cm} 
  	\end{beamercolorbox}
  %la deuxième ligne
	\begin{beamercolorbox}[ht=0.4cm, vmode]{subsection in head/foot}
	%titre de la section si elle est pas 0
		\ifnum\value{section}=0{} 
		\else{ \hspace{0.8cm} \thesection - \insertsectionhead }
		\fi
	%séparateur + titre de la sous-section si elle est pas 0
		\ifnum\value{subsection}=0{} 
		\else{ 
			\hspace*{0.1cm} \couleur{$\bullet$} \hspace*{0.1cm} 
			\thesection.\thesubsection \, \insertsubsectionhead
		}
		\fi
		\vspace*{0.12cm}
\end{beamercolorbox}
%\vspace*{-0.03cm} %pour pas qu'il y ait d'espace avec la ligne de frametitle
 }
%\setbeamertemplate{frametitleheigth}{4cm}
\setbeamertemplate{frametitle}{
	\vspace*{-0.04cm} 
	\begin{beamercolorbox}[ht=0.8cm,wd=\paperwidth, vmode]{frametitle}
		\hspace{0.3cm} \insertframetitle \vspace*{0.1cm}
	\end{beamercolorbox}
}

%commande pour ajuster l'alignement vertical des titres sans lettres descendantes
%\newcommand{\esp}{\\[0.1cm]} %--version qui marche sans le package minted
\newcommand{\esp}{\\[-0.5cm]} %--version qui marche avec le package minted


%--------couleurs
\setbeamercolor{structure}{fg=turquoiseFonce!70!black} 

\setbeamercolor{block title}{fg=turquoiseFonce!70!black,bg=vertdEau}
\setbeamercolor{block body}{bg=vertdEau!10!white}

\setbeamercolor{block title alerted}{bg=vertdEau!85!white,fg=turquoiseFonce!80!black}
\setbeamercolor{block body alerted}{bg=vertdEau!8!white}
%\setbeamercolor{alerted text}{fg=red}

\setbeamercolor{block title example}{bg=vertdEau,fg=turquoiseFonce}
\setbeamercolor{block body example}{bg=vertdEau!10!white}
\setbeamercolor{example text}{fg=blue!20!turquoise}

%-------- TOC
\setbeamertemplate{section in toc}[sections numbered]

%-----------------------------------------------
%Plan qui s'affiche au début de chaque section %|
\AtBeginSection[]{                             %|
\begin{frame}[plain]                           %|
\frametitle{Plan\\[0.1cm]}                     %|
\tableofcontents[                              %|
currentsection,                                %|
hideothersubsections,                          %|
subsubsectionstyle=hide]                       %|
\addtocounter{framenumber}{-1}                 %|
\end{frame}}                                   %|
%-----------------------------------------------




%-------- commande pour les ref sur les slides
\newcommand{\bandeauREF}[1]{
\noindent\makebox[\textwidth][l]{%
\hspace{-\dimexpr\oddsidemargin+1in}%
\colorbox{expli!20!white}{%
\parbox{\dimexpr\paperwidth-2\fboxsep}{
\footnotesize\textcolor{expli!80!black}{#1}
}}}}



%%%%%%%%%%%%%% MA PALETTE DE COULEURS
%%%%%%%%%%%%%% couleurs locales
\newcommand{\cearly}{orange!50!oranger}
\newcommand{\ctardy}{orange!50!oranger!50!magenta!50!violet}
\newcommand{\cmilieu}{yellow!45!magenta}
\newcommand{\cfenetre}{yellow!45!magenta}
%\newcommand{\cfenetre}{orange!80!magenta!80!violet}
\newcommand{\caxe}{blue!60!black}

\definecolor{vert}{RGB}{194, 247, 50}
\definecolor{vertClair}{RGB}{80, 180, 33}
\definecolor{vertFonce}{RGB}{20, 148, 20} 


\definecolor{turquoiseFonce}{RGB}{0, 149, 182} 
\definecolor{turquoise}{RGB}{0, 204, 203}
%\definecolor{turquoise}{RGB}{53, 123, 153} 
\colorlet{monCyan}{turquoise!90!blue} 

\colorlet{vertdEau}{blue!50!green!30!white} 

\definecolor{prec}{RGB}{128, 0, 128} 
\definecolor{hyp}{RGB}{255, 0, 127} 


\definecolor{orange}{RGB}{250, 130, 1}
\definecolor{orange1}{RGB}{244, 142, 3}
\definecolor{orange2}{RGB}{255, 195, 0}
%\colorlet{orangé}{orange!55!yellow}
\colorlet{oranger}{orange!55!yellow}
\definecolor{jaune}{RGB}{255, 195, 0 }
\colorlet{ocre}{orange!90!green!80!white} 


\definecolor{briqueRouge}{RGB}{200, 48, 24}
\definecolor{brique}{RGB}{141, 48, 24}
\definecolor{brown}{RGB}{165, 137, 107}
\definecolor{chamois}{RGB}{200, 141, 75}

\colorlet{lavande}{vertdEau!50!blue}
\definecolor{lilas}{RGB}{154, 107, 165}
\colorlet{parme}{magenta} %violet!30!magenta
\definecolor{mauve}{RGB}{161, 132, 220}%{147, 112, 219}

\definecolor{prunelle}{RGB}{106,24, 141}
\definecolor{prune}{RGB}{121, 7, 123}
\definecolor{violet}{RGB}{153, 0, 139}
\definecolor{violette}{RGB}{128, 0, 128} 

\definecolor{rose}{RGB}{255, 0, 127} 
\definecolor{framboise}{RGB}{141, 24, 59}
\colorlet{grenat}{magenta!90!red!80!black}

\colorlet{expli}{gray}

\definecolor{hellseahorse}{RGB}{204, 204, 255}
\definecolor{seahorse}{RGB}{204,180, 255}
\definecolor{darkseahorse}{RGB}{83, 74, 196}

%symboles pratiques
\newcommand{\flch}{\item[$\rightarrow$]}
\newcommand{\dc}{{\usebeamercolor[fg]{structure}$\hookrightarrow$}}
\newcommand{\ok}{\textcolor{vert}{\checkmark}}
\newcommand{\point}{{\usebeamercolor[fg]{structure}$\bullet\enskip$}}
\newcommand{\Point}{{\usebeamercolor[fg]{structure}$\bullet\enskip$}}


%styles
\newcommand{\couleur}[1]{{\usebeamercolor[fg]{structure}#1}}
\newcommand{\important}[1]{\couleur{\textbf{#1}}}
\newcommand{\remarque}[1]{\textit{\textrm{#1}}}

%pour le template
\newcommand{\lin}[1]{\mintinline{latex}{#1}}
%%%%%%%%%%%%%%%%%%%%%%%%%%%%%%%%%%%%%%%%%%%%%%%%%
\begin{document}
%

%%%%%%%%%%%%%%%%%%%%%%%%%%%%%%%%%%%%%%%%%%%%%%%%%
% A COMPLETER 
%%%%%%%%%%%%%%%%%%%%%%%%%%%%%%%%%%%%%%%%%%%%%%%%%

%entre crochet le titre court pour l'en tête puis le vrai titre
\title[La Transpilation]{
La transpilation de Fortran vers C
}

\author{
\large 
Présentation de \important{Yann MIQUEL--ERDMANN}\\[0.2cm]
Candidat n° \important{51280}\\[0.4cm]
\footnotesize
travail réalisé avec \couleur{Erwan FALAUX--BACHELOT}\\[0.4cm]
\vspace*{-1.2cm}
} 

%entre crochet la short date pour l'en-tête
\date[Juillet 2025]{}

\institute{}


%------------------------------------------------
\begin{frame}[plain]
\titlepage 
\addtocounter{framenumber}{-1} 
\end{frame}

%------------------------------------------------
\title[La transpilation de Fortran vers C]{}


\begin{frame}
<<<<<<< Updated upstream
<<<<<<< Updated upstream
    \frametitle{Introduction \esp}
=======
    \frametitle{Introduction\esp}

	Le Fortran: 
    \includegraphics[scale=0.03]{static/Fortran_logo.svg.png}\\
    Créé par: IMB



\end{frame}


\begin{frame}
    \frametitle{Problématique\esp}
>>>>>>> Stashed changes
    
	\vspace*{0.5cm}
=======
    \frametitle{Introduction\esp}
>>>>>>> Stashed changes

    \begin{minipage}[t]{0.48\textwidth}
        \includegraphics[scale=0.05]{static/Fortran_logo.svg.png}\\
        \textbf{Le Fortran:}\\ 
        Créé par:  John Backus, IBM\\
        Créé en:   1957\\
        Dernière version en: 2023\\
        Rang sur Github: $9^{\text{ème}}$
    \end{minipage}
    \hspace*{0.2cm}
    \begin{minipage}[t]{0.48\textwidth}
        \includegraphics[scale=0.05]{static/C_logo.png}\\
        \textbf{Le C:}\\
        Créé par: Dennis Ritchie et \\Kenneth Thompson, Bell Labs\\
        Créé en:  1972\\
        Dernière version en: 2024\\
        Rang sur Github: $41^{\text{ème}}$


    \end{minipage}
    
	\vspace*{2.9cm}

    \bandeauREF{%
        \hspace*{0.2cm}\textbf{Logo Fortran}: \url{https://commons.wikimedia.org/wiki/File:Fortran_logo.svg}\\
        \hspace*{0.2cm}\textbf{Logo C}: \url{https://commons.wikimedia.org/wiki/File:C_Logo.png}
    }

\end{frame}


\begin{frame}
    \frametitle{Problématique\esp}
    
    \begin{center}
        \textbf{Comment implémenter une conversion rapide de programmes Fortran en programmes C?}
    \end{center}

\end{frame}


\section{La transpilation}

\begin{frame}
    \frametitle{Définition}
    \begin{tikzpicture}
    
    \tikzstyle{lien} = [->, >=latex]
    \tikzstyle{basic_text}=[text width=2cm, text badly centered]
    \tikzstyle{basic_node}=[draw = black,rounded corners=4pt, basic_text]
    \tikzstyle{wrapper}=[basic_node, inner sep=3pt]
    \tikzstyle{hidden}=[draw=black!0,color=black!0]
    \tikzstyle{faded}=[draw=black!20, color=black!20]
    

    \node [basic_node, text width=5cm, align=left] (Fortran) {
        \scalebox{0.8}{
            \parbox{\textwidth}{
                \textbf{Programe Fortran}
                \inputminted[firstline=1, lastline=3,firstnumber=1]{fortran}{static/HelloWorld.f90}
            }
        }
    };

    \node [basic_node, text width=5cm, align=left] (C) [right=1cm of Fortran] {
            \scalebox{0.8}{
            \parbox{\textwidth}{
                \textbf{Programe C}

        \inputminted[firstline=1, lastline=4,firstnumber=1]{c}{static/HelloWorld.c}
            } } 
    };
 
    \draw [lien] (Fortran.north) to[out=45, in=135]  node[midway, above] {transpilation}   (C.north west) ; 

\end{tikzpicture}
\end{frame}


\begin{frame}

    \frametitle{Les étapes de la transpilation}

   \begin{center}

  \begin{tikzpicture}
    \tikzstyle{lien} = [->, >=latex]
    \tikzstyle{basic_text}=[text width=2cm, text badly centered, , minimum height=1cm]
    \tikzstyle{basic_node}=[draw = black,rounded corners=4pt, basic_text ]
    \tikzstyle{wrapper}=[basic_node, inner sep=3pt]

    
    \draw[white](-1.5,-4)rectangle(10,3);
      



    \node[basic_node] (A) at (0, 0) {analyse lexicale};
    \node[basic_node] (B) [right=0.5cm of A] {analyse syntaxique};
    \node[basic_node] (C) [right=0.5cm of B] {abstraction};

    \node[basic_node] (D) [right = 0.75cm of C] {traduction vers C};
    
    \draw [lien] (A) -- (B);
    \draw [lien] (B) -- (C);
    \draw [lien] (C) -- (D);
    
    
    \only{
      \node[basic_node] (E) [above = 0.75cm of D] {traduction vers Fortran};
      \node[basic_node] (F) [below = 0.75cm of D] {traduction vers Python};
      
      \node [wrapper, fit=(E), draw=red] (red_fit_2) {};
      \node [wrapper, fit=(F), draw=red] (red_fit_3) {};
      \draw [lien] (C) to [bend left] (E.west); 
      \draw [lien] (C) to [bend right] (F.west); 
    }<4->


    % \node[fit=(red_fit_1)(red_fit_2)(red_fit_3)(blue_fit)] (global_fit) {};
    
    \node<2-> [wrapper, fit=(A)(B)(C), draw=vertFonce] (blue_fit) {};
    \node<2-> [below=2cm of A.west, anchor=west, color=vertFonce] (label1) {module du langage d'entrée};
    
    \node<3-> [wrapper, fit=(D), draw=red] (red_fit_1) {};
    \node<3-> [below=0.5cm of label1.west, anchor=west, color=red] {modules des langages de sortie};

    
  \end{tikzpicture}
\end{center}
\end{frame}



\section{Analyse}

\subsection{Grammaire}
\begin{frame}
    \frametitle{Définition \esp}

    \textbf{Grammaire non contextuelle:}\\
    $\bullet \ G=(\Sigma, V, P, S)$\\
    $\bullet$ les règles de productions sont de la forme $W \rightarrow U_1...U_k\ et\ \forall i \in \llbracket1,\ k\rrbracket, U_i \in V \cup \Sigma , \ W \in V $
    \vspace{0.5cm}
    $ $\\
    {
        $\Sigma = \{a,b\}, V = \{A,B\}, S = A, P = \{$\\
        $ S \rightarrow A ;$\\
        $ A \rightarrow aA ;$\\
        $ A \rightarrow B ;$\\
        $ B \rightarrow bB ;$\\
        $ B \rightarrow \epsilon$\\
        $\}$\\
        $\mathcal{L} (G) = a^*b^*$
    }
\end{frame}

\begin{frame}
    \frametitle{Un exemple de grammaire\esp}
    \inputminted[breaklines, fontsize=\small]{text}{static/Grammaire_exemple.txt}
\end{frame}

\begin{frame}
    \frametitle{Nos types ocaml\esp}
    \noindent
    \begin{minipage}[t]{0.48\textwidth}
        \scalebox{0.9}{
            \parbox{\textwidth}{
                \inputminted[breaklines, fontsize=\smaller]{ocaml}{static/type_symbol.ml}
            }
        }
    \end{minipage}
        \hfill
    \begin{minipage}[t]{0.48\textwidth}
        \scalebox{0.9}{
        \parbox{\textwidth}{
        \inputminted[breaklines, fontsize=\small]{ocaml}{static/type_grammaire.ml}        
        }
        }
    \end{minipage}

\end{frame}




\subsection{Analyse lexicale}

\begin{frame}
    
   \input{tikz/Fortran_vers_Lexèmes.tex}

\end{frame}

\begin{frame}
  \frametitle{Définition d'un automate\esp}

  
  Les automates sont définis par $(\mathcal{A} = \Sigma, Q, I, F, \delta)$\\

  \vspace{0.5cm}
  \pause
  \begin{tikzpicture}
    \node [state, initial]         (1) at (0, 0)        {1};
    \node [state]                  (2) [right=1cm of 1] {2};
    \node [state, accepting right] (3) [right=1cm of 2] {3};

    \path [->]
    (1) edge              node [above] {$a$} (2)
    (2) edge              node [above] {$c$} (3)
    (2) edge [loop above] node [above] {$b$} (2)
    ;

  \end{tikzpicture}\\
  $\mathcal{L} (\mathcal{A} ) = a\cdot b^*\cdot c$

  
\end{frame}

%------------étude d'un exemple-----------
\begin{frame}
  \frametitle{Les automates: étude d'un exemple \esp}
    \begin{tikzpicture}
    \node at (3, 0) (text) {\ttfamily print*,"Hello World"};
    \begin{scope}
      [every node/.style={scale=0.6,font=\Large}, scale=0.4, initial text=, shift={(-13, 0)}]
      \tikzstyle{style0}=[]
      \tikzstyle{style1}=[]
      \tikzstyle{style2}=[]
      \tikzstyle{style3}=[]
      \tikzstyle{style4}=[] 
      \tikzstyle{style5}=[]
      \tikzstyle{style6}=[]
      \tikzstyle{style7}=[]
      \tikzstyle{style8}=[]
      \tikzstyle{style9}=[]

      \only<1> {\tikzstyle{style0}=[draw=blue!50,very thick,fill=blue!20]}
      \only<2> {\tikzstyle{style1}=[draw=blue!50,very thick,fill=blue!20]}
      \only<3> {\tikzstyle{style2}=[draw=blue!50,very thick,fill=blue!20]}
      \only<4> {\tikzstyle{style3}=[draw=blue!50,very thick,fill=blue!20]}
      \only<5> {\tikzstyle{style4}=[draw=blue!50,very thick,fill=blue!20]}
      \only<6> {\tikzstyle{style5}=[draw=blue!50,very thick,fill=blue!20]}
      \only<7> {\tikzstyle{style0}=[draw=blue!50,very thick,fill=blue!20]}
      \only<8> {\tikzstyle{style8}=[draw=blue!50,very thick,fill=blue!20]}
      \only<9> {\tikzstyle{style0}=[draw=blue!50,very thick,fill=blue!20]}
      \only<10>{\tikzstyle{style9}=[draw=blue!50,very thick,fill=blue!20]}
      \only<11> {\tikzstyle{style0}=[draw=blue!50,very thick,fill=blue!20]}
      \only<12-23>{\tikzstyle{style6}=[draw=blue!50,very thick,fill=blue!20]}
      \only<24>{\tikzstyle{style7}=[draw=blue!50,very thick,fill=blue!20]}

      \node[state, style0, initial        ]  at (0,0) (0)          {$I$};

      \node[state, style1] (1) at (2,1) {1};
      \node[state, style2] (2) [right=0.4cm of 1] {2};
      \node[state, style3] (3) [right=0.4cm of 2] {3};
      \node[state, style4] (4) [right=0.4cm of 3] {4};
      \node[state, accepting right, style5] (5) [right=0.4cm of 4] {$F_1$};
      
      \node[state, style6,                ] (6) at (2,3) {5};
      \node[state, style7, accepting right] (7) [right=0.4cm of 6] {$F_2$};
      
      \node[state, style8, accepting right] (8) at (2,-1)      {$F_3$};
      \node[state, style9, accepting right] (9) at (2,-3)      {$F_4$};

      \path[->]
      (0) edge                node[above] {$*$} (8)
      (0) edge                node[above] {\LARGE$,$} (9)

      (0) edge                node[above] {$p$} (1)
      (1) edge                node[above] {$r$} (2)
      (2) edge                node[above] {$i$} (3)
      (3) edge                node[above] {$n$} (4)
      (4) edge                node[above] {$t$} (5)
      
      
      (0) edge                node[above] {$"$} (6)
      (6) edge [loop above]   node[above] {$[a-z]|[A-Z]$} (6)
      (6) edge                node[above] {$"$} (7)
      ;
      
    \end{scope}
    \foreach \step in {1,2,...,6}{
      \node<\step>[right=\fpeval{\step * 5.75 - 3}pt of text.west, draw=blue, fill=blue!20, inner sep=0pt, minimum width=2pt, minimum height=10pt]{};
    }
    \foreach \step in {7,8}{
      \node<\step>[right=\fpeval{(\step - 1) * 5.75 - 3}pt of text.west, draw=blue, fill=blue!20, inner sep=0pt, minimum width=2pt, minimum height=10pt]{};
    }
    \foreach \step in {9,10}{
      \node<\step>[right=\fpeval{(\step - 2) * 5.75 - 3}pt of text.west, draw=blue, fill=blue!20, inner sep=0pt, minimum width=2pt, minimum height=10pt]{};
    }
    \foreach \step in {11,12,...,24}{
      \node<\step>[right=\fpeval{(\step - 3) * 5.75 - 3}pt of text.west, draw=blue, fill=blue!20, inner sep=0pt, minimum width=2pt, minimum height=10pt]{};
    }

    \node at (0, -2.3) (lb) {
      \begin{minipage}{\textwidth}
        \only<1-6>  {\inputminted[firstline=1, lastline=1, linenos=false, frame=none, xleftmargin=0pt,]{ocaml}{static/exampleList.ml}}
        \only<7-8>  {\inputminted[firstline=2, lastline=2, linenos=false, frame=none, xleftmargin=0pt,]{ocaml}{static/exampleList.ml}}
        \only<9-10> {\inputminted[firstline=3, lastline=3, linenos=false, frame=none, xleftmargin=0pt,]{ocaml}{static/exampleList.ml}}
        \only<11-24>{\inputminted[firstline=4, lastline=4, linenos=false, frame=none, xleftmargin=0pt,]{ocaml}{static/exampleList.ml}}
        \only<25>   {\inputminted[firstline=5, lastline=5, linenos=false, frame=none, xleftmargin=0pt,]{ocaml}{static/exampleList.ml}}
      \end{minipage}
    };
  \end{tikzpicture}
\end{frame}


%------------étude d'un exemple-----------
\begin{frame}
  \frametitle{Les automates: étude d'un exemple\esp}
  \begin{tikzpicture}
    \node at (1, 0) (text1) {\ttfamily print};
    \node [right=-5pt of text1.east, anchor=west] (text2) {\ttfamily print};
    \node [right=-5pt of text2.east, anchor=west] (text3) {\ttfamily "Hello World"};

    \begin{scope}
      [every node/.style={scale=0.6,font=\Large}, scale=0.4, initial text=, shift={(-13, 0)}]
      \tikzstyle{style0}=[]
      \tikzstyle{style1}=[]
      \tikzstyle{style2}=[]

      \only<1> {\tikzstyle{style0}=[draw=blue!50,very thick,fill=blue!20]}
      \only<2-3>{\tikzstyle{style1}=[draw=blue!50,very thick,fill=blue!20]}
      \only<4> {\tikzstyle{style2}=[draw=blue!50,very thick,fill=blue!20]}

      \node[state, style0, style1, style2, initial]  at (0,0) (0)          {$I$};

      \node[state, style1] (1) at (2,1) {1};
      \node[state, style1] (2) [right=0.4cm of 1] {2};
      \node[state, style1] (3) [right=0.4cm of 2] {3};
      \node[state, style1] (4) [right=0.4cm of 3] {4};
      \node[state, accepting right, style1] (5) [right=0.4cm of 4] {$F_1$};
      
      \node[state, style2,                ] (6) at (2,3) {5};
      \node[state, style2, accepting right] (7) [right=0.4cm of 6] {$F_2$};
      
      \node[state, accepting right]         (8) at (2,-1)      {$F_3$};
      \node[state, accepting right]         (9) at (2,-3)      {$F_4$};

      \path[->]
      (0) edge                      node[above] {$*$} (8)
      (0) edge                      node[above] {\LARGE$,$} (9)

      (0) edge [style1]             node[above] {$p$} (1)
      (1) edge [style1]             node[above] {$r$} (2)
      (2) edge [style1]             node[above] {$i$} (3)
      (3) edge [style1]             node[above] {$n$} (4)
      (4) edge [style1]             node[above] {$t$} (5)
      
      
      (0) edge [style2]             node[above] {$"$} (6)
      (6) edge [loop above, style2] node[above] {$[a-z]|[A-Z]$} (6)
      (6) edge [style2]             node[above] {$"$} (7)
      ;
      
    \end{scope}
    
    \node<2>[fit=(text1), draw=blue, fill=blue!20, inner sep=-2pt, fill opacity=0.5]{};
    \node<3>[fit=(text2), draw=blue, fill=blue!20, inner sep=-2pt, fill opacity=0.5]{};
    \node<4>[fit=(text3), draw=blue, fill=blue!20, inner sep=-2pt, fill opacity=0.5]{};
    \node at (0, -2.3) {
      \begin{minipage}{\textwidth}
        \only<1>{\inputminted[firstline=1, lastline=1, linenos=false, frame=none, xleftmargin=0pt,]{ocaml}{static/exampleList.ml}}
        \only<2>{\inputminted[firstline=2, lastline=2, linenos=false, frame=none, xleftmargin=0pt,]{ocaml}{static/exampleList.ml}}
        \only<3>{\inputminted[firstline=6, lastline=6, linenos=false, frame=none, xleftmargin=0pt,]{ocaml}{static/exampleList.ml}}
        \only<4>{\inputminted[firstline=7, lastline=7, linenos=false, frame=none, xleftmargin=0pt,]{ocaml}{static/exampleList.ml}}
      \end{minipage}
    };
  \end{tikzpicture}
\end{frame}




\subsection{Analyse Syntaxique}

\begin{frame}
    \input{tikz/Lexèmes_vers_Arbre_Syntaxique.tex}
\end{frame}

\begin{frame}
    \frametitle{Algorithme LL1\esp}

    Nécessite:\\
    \vspace{0.05cm}\\{
    $\bullet$ Une grammaire LL (sans récursion gauche et sans ambigüité)\\
    $\text{\textcolor{red}{\ding{55}}} \ \ \mathcal{X} \rightarrow \mathcal{X}a$\\
    $\text{\textcolor{red}{\ding{55}}} \ \  \mathcal{X} \rightarrow a\mathcal{X}\mathcal{X} \ | \ \varepsilon  $\\}
    \tikzstyle{feuille}=[draw, rectangle, inner sep = 0.12cm]
\tikzstyle{noeud}=[draw, rectangle,rounded corners, minimum width= 0.64cm, line width = 1pt]
\begin{tikzpicture}[
        baseline=(base), 
        level/.style={sibling distance = 1cm/#1, level distance = 0.5cm},
        every node/.style={scale=0.7, font=\smaller}
    ]
    
    \node[ noeud] (1) {X}
    child { 
        node[feuille] {a}
    }
    child { 
        node [ noeud] {X}
        child {
            node[feuille] {a}
        }
    }
    child { 
        node[ noeud] {X}
        child { 
            node [feuille] {$\varepsilon$}
        }
    };

    \node[ noeud, right= 3cm of 1] {X}
    child { 
        node[feuille] {a}
    }
    child { 
        node[ noeud] {X}
        child { 
            node [feuille] {$\varepsilon$}
        }
    }
    child { 
        node [ noeud] {X}
        child {
            node[feuille] {a}
        }
    };

\end{tikzpicture}

    
    \vspace{0.25cm}
    $\bullet$ LL1: un seul lexème suffit à choisir la règle de production\\
    $\mathcal{X} \rightarrow a\mathcal{X} | b\mathcal{Y} $\\
    $\mathcal{Y} \rightarrow b\mathcal{Y} | d\mathcal{Y} | \varepsilon $
    
\end{frame}

%%%%%%%%%%%%%%%% pas pour Erwan
\begin{frame}
    \frametitle{Tableau des premiers et suivants\esp}
    \small
    (1): \textcolor{blue}{Programme} -> \textcolor{blue}{ProgramMC NomProg Print EndProgramMC NomProg}\\
    (2): \textcolor{blue}{Print} -> \textcolor{blue}{PrintMC Asterisque Virgule Chaine ParamListe}\\
    (3): \textcolor{blue}{ParamListe} -> \textcolor{blue}{Virgule Chaine ParamListe}\\
    (4): \textcolor{blue}{ParamListe} ->  \textcolor{blue}{$\varepsilon$} \\

    \vspace{0.25cm}{
    \begin{tabular}{|c|c|c|}
    \hline
        & Premiers & Suivants\\  
    \hline
    \textcolor{blue}{Programme} & $\{\text{\textcolor{blue}{ProgramMC}}\}$ & $\{\textcolor{blue}{EOF}\}$ \\
    \hline
    \textcolor{blue}{Print} & $\{\text{\textcolor{blue}{PrintMC}}\}$ & $\{\text{\textcolor{blue}{EndProgramMC}}\}$ \\
    \hline
    \textcolor{blue}{ParamListe} & $\{ \text{\textcolor{blue}{Virgule}}, \textcolor{blue}{\varepsilon} \}$ & $\{\text{\textcolor{blue}{EndProgramMC}}\}$\\
    \hline

    \end{tabular}
    }
    \vspace{0.25cm}
    {
    \pause
    
        \footnotesize
        \begin{tabular}{|c|c|c|c|c|c|c|c|c|c|}
            \hline
            & \textcolor{blue}{PrintMC} & \textcolor{blue}{ProgramMC} & \textcolor{blue}{EndProgramMC} & \textcolor{blue}{Virgule} \\
            \hline
            \textcolor{blue}{Programme} & & (1) & & \\
            \hline
            \textcolor{blue}{Print} & (2) & & & \\
            \hline
            \textcolor{blue}{ParamListe} & & & (4) & (3) \\
            \hline
        \end{tabular}
    
        }   
\end{frame}



\begin{frame}
    \frametitle{Arbre de syntaxe\esp}
    \begin{center}
        \scalebox{1.2}{
        \tikzstyle{feuille}=[draw, rectangle, inner sep = 0.12cm]
\tikzstyle{noeud}=[draw, rectangle,rounded corners, minimum width= 0.64cm, line width = 1pt]
\begin{tikzpicture}[
        baseline=(base), 
        level/.style={sibling distance = 2cm/#1, level distance = 1cm},
        every node/.style={scale=0.6, font=\footnotesize}
    ]
    
    \node[ noeud] {Programme}
    child { 
        node[ noeud] (base) {ProgramMC}
        child { 
            node [ feuille] {program}
        }
    }
    child { 
        node[ noeud] {NLigne}
        child { 
            node [ feuille] {\textbackslash n}
        }
    }
    child [ level distance=2cm]{ 
        node [ noeud] {Print}
        child {
            node [noeud] {PrintMC}
            child { 
                node [ feuille] {print}
            }
        }
        child{
            node [noeud] {Asterisque}
            child { 
                node [feuille] {*}
            }
        }
        child{
            node [noeud] {Virgule}
            child { 
                node [feuille] {,}
            }
        }
        child{
            node [noeud] {Chaine}
            child { 
                node [feuille] {"Hello World"}
            }
        }
        child{
            node [noeud] {Paramliste}
            child { 
                node [feuille] {$\varepsilon$}
            }
        }
    }
    child { 
        node[ noeud] {EndProgramMC}
        child { 
            node [ feuille] {end program}
        }
    };

\end{tikzpicture}
        }

    \end{center}
\end{frame}



\begin{frame}
    \frametitle{Conversion en arbre de syntaxe abstraite\esp}
    \begin{tikzpicture}
    
    \tikzstyle{lien} = [->, >=latex]
    \tikzstyle{basic_text}=[text width=2cm, text badly centered]
    \tikzstyle{basic_node}=[draw = black,rounded corners=4pt, basic_text]
    \tikzstyle{wrapper}=[basic_node, inner sep=3pt]
    \tikzstyle{hidden}=[draw=black!0,color=black!0]
    \tikzstyle{faded}=[draw=black!20, color=black!20]
    

    \node [basic_node, text width=6cm, align=left] (AS) {
         \scalebox{0.7}{
            \parbox{\textwidth}{
                \tikzstyle{feuille}=[draw, rectangle, inner sep = 0.12cm]
\tikzstyle{noeud}=[draw, rectangle,rounded corners, minimum width= 0.64cm, line width = 1pt]
\begin{tikzpicture}[
        baseline=(base), 
        level/.style={sibling distance = 2cm/#1, level distance = 1cm},
        every node/.style={scale=0.6, font=\footnotesize}
    ]
    
    \node[ noeud] {Programme}
    child { 
        node[ noeud] (base) {ProgramMC}
        child { 
            node [ feuille] {program}
        }
    }
    child { 
        node[ noeud] {NLigne}
        child { 
            node [ feuille] {\textbackslash n}
        }
    }
    child [ level distance=2cm]{ 
        node [ noeud] {Print}
        child {
            node [noeud] {PrintMC}
            child { 
                node [ feuille] {print}
            }
        }
        child{
            node [noeud] {Asterisque}
            child { 
                node [feuille] {*}
            }
        }
        child{
            node [noeud] {Virgule}
            child { 
                node [feuille] {,}
            }
        }
        child{
            node [noeud] {Chaine}
            child { 
                node [feuille] {"Hello World"}
            }
        }
        child{
            node [noeud] {Paramliste}
            child { 
                node [feuille] {$\varepsilon$}
            }
        }
    }
    child { 
        node[ noeud] {EndProgramMC}
        child { 
            node [ feuille] {end program}
        }
    };

\end{tikzpicture}
            }
        }
    };

    \node [basic_node, text width=4cm, align=left] (AST) [right=1cm of AS] {
        \scalebox{0.7}{
            \parbox{\textwidth}{
                \begin{tikzpicture}

    \node[font=\scriptsize] (0) at (0, 0) {Programme};
    \node[font=\scriptsize] (1) at (-1, -1) {NLigne};
    \node[font=\scriptsize] (2) at (1, -1) {Print};
    \node[font=\scriptsize] (3) at (1, -2) {'"Hello World"'};
    
    \draw (0) -- (1);
    \draw (0) -- (2);
    \draw (2) -- (3);
\end{tikzpicture}
            }
        }
    };
 
    \draw [lien] (AS.north) to[out=45, in=135]  node[midway, above] {\textbf{Conversion en arbre de syntaxe abstraite}}  (AST.north west) ; 

\end{tikzpicture}
\end{frame}



\section{Traduction}

\begin{frame}
    \begin{tikzpicture}
    
    \tikzstyle{lien} = [->, >=latex]
    \tikzstyle{basic_text}=[text width=2cm, text badly centered]
    \tikzstyle{basic_node}=[draw = black,rounded corners=4pt, basic_text]
    \tikzstyle{wrapper}=[basic_node, inner sep=3pt]
    \tikzstyle{hidden}=[draw=black!0,color=black!0]
    \tikzstyle{faded}=[draw=black!20, color=black!20]
    

    \node [basic_node, text width=4.7cm, align=left] (AST) {
         \scalebox{0.8}{
            \parbox{\textwidth}{
                \textbf{Arbre de syntaxe abstraite}
                \begin{tikzpicture}

    \node[font=\scriptsize] (0) at (0, 0) {Programme};
    \node[font=\scriptsize] (1) at (-1, -1) {NLigne};
    \node[font=\scriptsize] (2) at (1, -1) {Print};
    \node[font=\scriptsize] (3) at (1, -2) {'"Hello World"'};
    
    \draw (0) -- (1);
    \draw (0) -- (2);
    \draw (2) -- (3);
\end{tikzpicture}
            }
        }
    };

    \node [basic_node, text width=5.3cm, align=left] (C) [right=1cm of AST] {
        \scalebox{0.9}{
            \parbox{\textwidth}{
                \textbf{Programme C}

                \inputminted[firstline=1, lastline=4,firstnumber=1]{c}{static/HelloWorld.c}
            } 
        } 
    };
 
    \draw [lien] (AST.north) to[out=45, in=135]  node[midway, above] {\textbf{Traduction en C}}  (C.north west) ; 

\end{tikzpicture}
\end{frame}

\begin{frame}
    \frametitle{Traduction de L'arbre de syntaxe abstraite\esp}
    \begin{center}
        \begin{tikzpicture}

    \node[font=\scriptsize] (0) at (0, 0) {Programme};
    \node[font=\scriptsize] (1) at (-1, -1) {NLigne};
    \node[font=\scriptsize] (2) at (1, -1) {Print};
    \node[font=\scriptsize] (3) at (1, -2) {'"Hello World"'};
    
    \draw (0) -- (1);
    \draw (0) -- (2);
    \draw (2) -- (3);
\end{tikzpicture}
    \end{center}
    \vspace{0.5cm}
    On écrit le code C à l'aide d'un parcours en profondeur de l'AST


\end{frame}


\section{Notre transpileur}

\begin{frame}
    on a donc fait un exécutable qui peut transpiler du Fortran vers du C ou du Fortran
    il s'utilise de la manière suivante
    explication -C -Fortran et -o 

\end{frame}


\begin{frame}
    Tests, nous avons essayé notre transpileur sur plusieurs programmes tests en fortran tels que 
    donner des exemple d'algos
    et il fonctionne
\end{frame}


\section{Conclusion}


\begin{frame}
    \frametitle{Conclusion\esp}

    \point Conversion rapide\\
    \point Une partie du transpileur est construite automatiquement\\
    \point La syntaxe prise en charge est étendue\\
    \point Les parties du transpileur sont modulaires 

\end{frame}


%------------------------------------------------
% \tikzstyle{feuille}=[draw, rectangle, inner sep = 0.12cm]
\tikzstyle{noeud}=[draw, rectangle,rounded corners, minimum width= 0.64cm, line width = 1pt]
\begin{tikzpicture}[
        baseline=(base), 
        level/.style={sibling distance = 2cm/#1, level distance = 1cm},
        every node/.style={scale=0.6, font=\footnotesize}
    ]
    
    \node[ noeud] {Programme}
    child { 
        node[ noeud] (base) {ProgramMC}
        child { 
            node [ feuille] {program}
        }
    }
    child { 
        node[ noeud] {NLigne}
        child { 
            node [ feuille] {\textbackslash n}
        }
    }
    child [ level distance=2cm]{ 
        node [ noeud] {Print}
        child {
            node [noeud] {PrintMC}
            child { 
                node [ feuille] {print}
            }
        }
        child{
            node [noeud] {Asterisque}
            child { 
                node [feuille] {*}
            }
        }
        child{
            node [noeud] {Virgule}
            child { 
                node [feuille] {,}
            }
        }
        child{
            node [noeud] {Chaine}
            child { 
                node [feuille] {"Hello World"}
            }
        }
        child{
            node [noeud] {Paramliste}
            child { 
                node [feuille] {$\varepsilon$}
            }
        }
    }
    child { 
        node[ noeud] {EndProgramMC}
        child { 
            node [ feuille] {end program}
        }
    };

\end{tikzpicture}

% \section{Annexe}

% \tableofcontents

\subsection{\url{generate_ml_files.ml}}
\inputminted[]{ocaml}{../../src/preprocessing/generate_ml_files.ml}
\pagebreak
\subsection{\url{generate_grammar.ml}}
\inputminted[]{ocaml}{../../src/preprocessing/generate_grammar.ml}
\pagebreak
\subsection{\url{build_automaton.ml}}
\inputminted[]{ocaml}{../../src/prebuild/build_automaton.ml}
\pagebreak
\subsection{\url{abstract_tokens.ml}}
\inputminted[]{ocaml}{../../src/abstract_tokens.ml}
\pagebreak
\subsection{\url{bibliotheques.ml}}
\inputminted[]{ocaml}{../../src/bibliotheques.ml}
\pagebreak
\subsection{\url{environnement.ml}}
\inputminted[]{ocaml}{../../src/environnement.ml}
\pagebreak
\subsection{\url{symbols.ml}}
\inputminted[]{ocaml}{../../src/symbols.ml}
\pagebreak
\subsection{\url{traduction.ml}}
\inputminted[]{ocaml}{../../src/traduction.ml}
\pagebreak
\subsection{\url{regex.ml}}
\inputminted[]{ocaml}{../../src/regex.ml}
\pagebreak
\subsection{\url{grammar_functions.ml}}
\inputminted[]{ocaml}{../../src/grammar_functions.ml}
\pagebreak
\subsection{\url{grammar.ml}}
\inputminted[]{ocaml}{../../src/grammar.ml}
\pagebreak
\subsection{\url{automates.ml}}
\inputminted[]{ocaml}{../../src/automates.ml}
\pagebreak
\subsection{\url{LL1.ml}}
\inputminted[]{ocaml}{../../src/LL1.ml}
\pagebreak
\subsection{\url{convert_to_abstract.ml}}
\inputminted[]{ocaml}{../../src/convert_to_abstract.ml}
\pagebreak
\subsection{\url{det_automaton.ml}}
\inputminted[]{ocaml}{../../src/det_automaton.ml}
\pagebreak
\subsection{\url{traductionC.ml}}
\inputminted[]{ocaml}{../../src/traductionC.ml}
\pagebreak
\subsection{\url{traductionFortran.ml}}
\inputminted[]{ocaml}{../../src/traductionFortran.ml}
\pagebreak
\subsection{\url{transpileurs.ml}}
\inputminted[]{ocaml}{../../src/transpileurs.ml}
\pagebreak
\subsection{\url{main.ml}}
\inputminted[]{ocaml}{../../src/main.ml}



\end{document}
